The router uses a deterministic XY routing algorithm in the 3×3 mesh network. The output port is selected by comparing the current router coordinates with the destination coordinates in the header flit.

The packet is first routed in the X direction. If the destination is on the right, the packet is sent to EAST; if it is on the left, it is sent to WEST. When the X coordinate is already matched, the router then compares the Y coordinate and sends the packet to NORTH or SOUTH. If both coordinates match, the packet is delivered to the LOCAL port.

This method is simple and hardware-efficient, since it only requires coordinate comparison and no routing table. Another advantage is that XY routing naturally restricts all turns from the Y dimension back to the X dimension. Therefore, cyclic channel dependencies are avoided, which helps prevent routing deadlock. Its limitation is that the path is fixed and cannot adapt to congestion.

\begin{algorithm}[htbp]
\caption{Pseudo code of the XY-routing computation}
\label{alg:xy_routing}
\textbf{Definitions:}

\begin{itemize}
    \item $cur_x, cur_y$: current router coordinates.
    \item $dst_x, dst_y$: destination router coordinates.
    \item $port$: selected output port.
\end{itemize}

\textbf{Input:}

\begin{itemize}
    \item $cur_x, cur_y$: coordinates of current router.
    \item $dst_x, dst_y$: coordinates extracted from the header flit.
\end{itemize}

\textbf{Output:}

\begin{itemize}
    \item $port$: routing decision toward the next hop.
\end{itemize}

\begin{algorithmic}[1]
\If{$dst_x > cur_x$}
    \State $port \gets EAST$
\ElsIf{$dst_x < cur_x$}
    \State $port \gets WEST$
\ElsIf{$dst_y > cur_y$}
    \State $port \gets NORTH$
\ElsIf{$dst_y < cur_y$}
    \State $port \gets SOUTH$
\Else
    \State $port \gets LOCAL$
\EndIf
\State \Return $port$
\end{algorithmic}
\end{algorithm}

